% !TeX root = ../../Book1.tex
% 纲要标题：### 9.3 结构定理
% 纲要位置：计划纲要.md，第 370 行。

\section{结构定理}
\label{b1:sec:0903}

% 纲要内容（仅作写作计划，尚非教材正文）：
%
% * 不变因子分解与初等因子分解
% * 结构定理的初等证明
% * 分解的唯一性
%

\subsection{不变因子分解与初等因子分解}
\label{b1:subsec:0903:01}

本节的目标是把任意有限生成阿贝尔群化为循环群的有限直和，而且给出可由群本身恢复的分类数据。有限循环因子有两种常用排列方式：按整除关系排列，或按素数幂拆开。

\begin{theorem}\label{b1:thm:fga-structure}
每个有限生成阿贝尔群 \(A\) 都有分解
\[
A\cong\mathbb Z^r\oplus C_{d_1}\oplus\cdots\oplus C_{d_s},
\qquad 1<d_1\mid d_2\mid\cdots\mid d_s,
\]
其中 \(r,s\ge0\)。整数 \(r\) 及序列 \(d_1,\ldots,d_s\) 由 \(A\) 唯一确定。也可以唯一地写为一个自由部分与若干素数幂阶循环群的直和，唯一性允许重新排列这些循环因子。
\end{theorem}
第一种称为\term[bubian-yinzi-fenjie]{不变因子分解}[invariant factor decomposition]，其中 \(d_i\) 为不变因子；第二种称为\term[chudeng-yinzi-fenjie]{初等因子分解}[elementary divisor decomposition]，各循环因子的素数幂阶称为初等因子。下文依次证明存在性与唯一性。

\ProseParagraphBreak
例如 \(C_4\oplus C_6\) 还没有按整除关系排列，因为 \(4\nmid6\)。先按素数拆开，
\[
C_4\oplus C_6\cong C_4\oplus C_2\oplus C_3
\cong C_2\oplus C_{12}.
\]
初等因子是 \(4,2,3\)，不变因子是 \(2,12\)。这两种写法表达同一个群，但各自必须遵守相应的归一化条件；不能在任意循环因子拆分之间直接谈唯一性。

\subsection{结构定理的初等证明}
\label{b1:subsec:0903:02}

生成元与关系构成一个整数矩阵。要整理关系，只能用可逆的整数行、列操作，不能像实数线性代数那样任意除以主元；除法会改变整数格。

\begin{lemma}[整数矩阵的 Smith 化简]\label{b1:lem:smith-integer}
对任意 \(m\times n\) 整数矩阵 \(M\)，可以通过交换两行或两列、把一行或一列乘 \(-1\)、把某行或列的整数倍加到另一行或列，化为矩形对角阵
\[
\operatorname{diag}(d_1,\ldots,d_t,0,\ldots),
\qquad 0<d_1\mid d_2\mid\cdots\mid d_t.
\]
\end{lemma}
\begin{proof}
若矩阵全零则结束。否则把一个非零元素移到左上角，并使其为正，记为 \(a\)。对第一列的某个元素 \(b\)，作整数除法 \(b=qa+r\)，其中 \(0\le r<a\)，再把第一行的 \(q\) 倍从该行减去。如果 \(r>0\)，交换该行与第一行，使正主元严格变小；若 \(r=0\)，就消去了该元素。对第一行作同样的列操作。每当出现非零余数，主元严格下降，因此这种下降只可能发生有限次；主元不下降时，有限步内可把第一行和第一列的其余元素都清零。

此时矩阵形如
\[
\begin{pmatrix}a&0\\0&B\end{pmatrix}.
\]
如果 \(a\) 不整除 \(B\) 中的某个元素 \(b_{ij}\)，把该元素所在的行加到第一行。左上角仍为 \(a\)，而第一行现在含有 \(b_{ij}\)。再次用列的整数除法就产生小于 \(a\) 的非零余数，使主元严格下降，然后重新清理首行首列。正整数不能无限下降，所以最终获得这样的分块形式，且 \(a\) 整除 \(B\) 的全部元素。

固定首行首列，对 \(B\) 递归进行相同步骤。行列数都减少，递归必停。由于 \(B\) 的所有元素一开始均为 \(a\) 的倍数，后续整数行列操作保持此性质，所以它所得的非零对角元也都是 \(a\) 的倍数。再用递归中的整除链，便得到全部 \(d_i\mid d_{i+1}\)。所用每一步操作都有整数逆操作，没有改变相应整数格的可逆坐标性质。
\end{proof}
上述结果称为整数矩阵的\term[smith-biaozhunxing]{Smith 标准形}[Smith normal form]。证明给出的步骤是一个终止算法，而不只是允许“选取适当基”的存在性叙述。

\begin{proof}[结构定理的存在性证明]
取 \(A\) 的有限生成元，得到满同态 \(\pi:\mathbb Z^m\to A\)。由\cref{b1:thm:free-abelian-subgroup}，核 \(K\) 有有限基 \(k_1,\ldots,k_t\)，其中 \(t\le m\)。把它们在标准基下的坐标作为列，组成 \(m\times t\) 整数矩阵 \(M\)。

整数可逆列操作只改变 \(K\) 的生成基，因而不改变其生成子群；整数可逆行操作是母群 \(\mathbb Z^m\) 的自同构，因而把商群变成同构的商群。对 \(M\) 作 Smith 化简后，核变成由
\[
d_1e_1,\ldots,d_te_t
\]
生成的子群。这里恰有 \(t\) 个非零对角元：原列在整数上无关，也就在有理数上无关，因为有理数关系乘以公共分母就成为整数关系；可逆行列操作保持这个性质。于是
\[
A\cong
\mathbb Z^m/K
\cong C_{d_1}\oplus\cdots\oplus C_{d_t}\oplus\mathbb Z^{m-t}.
\]
等于一的对角元只贡献平凡群，删去它们便得到不变因子形式。

最后把每个 \(d_i\) 作素因数分解。对互素的 \(u,v\)，映射
\(C_{uv}\to C_u\oplus C_v\)、\(\bar x\mapsto(\bar x,\bar x)\)
的核平凡，且两边元素数相同，所以为同构。反复应用便把每个有限循环因子拆成素数幂阶循环群，得到初等因子形式。
\end{proof}
证明同时得到：有限生成无挠阿贝尔群自由；一般有限生成阿贝尔群的挠子群可以分离成一个直和因子。不过具体选出的自由补群依赖基的选择，结构定理没有宣称这个补群唯一。

\subsection{分解的唯一性}
\label{b1:subsec:0903:03}

先恢复自由部分，再分别恢复每个素数的幂阶。挠子群 \(T=A_{\mathrm{tor}}\) 是由群内条件确定的，所以同构必把它映到另一个群的挠子群。上面的分解说明 \(A/T\cong\mathbb Z^r\)，于是\cref{b1:prop:free-abelian-rank}确定了 \(r\)。

\ProseParagraphBreak
对有限群 \(T\) 的 \(p\) 主分量 \(B=T_p\)，考虑全由整数倍定义的子群及商群
\[
B\supseteq pB\supseteq p^2B\supseteq\cdots,\qquad
p^{j-1}B/p^jB.
\]
若 \(B\cong\bigoplus_{i=1}^k C_{p^{e_i}}\)，每个因子对第 \(j\) 层的贡献在 \(e_i\ge j\) 时有 \(p\) 个元素，在 \(e_i<j\) 时只有一个。因此
\[
\bigl|p^{j-1}B/p^jB\bigr|=p^{c_j},
\qquad c_j=\#\{\,i\mid e_i\ge j\,\}.
\]
左边是由 \(B\) 本身确定的数，所以每个 \(c_j\) 都是同构不变量。恰好等于 \(j\) 的指数出现 \(c_j-c_{j+1}\) 次，因而全部 \(e_i\) 及其重数唯一。这证明初等因子的唯一性，没有引入模或张量积。

\ProseParagraphBreak
由初等因子恢复不变因子也没有歧义。若 \(T=0\)，有限循环因子为空，取 \(s=0\) 即可。以下设 \(T\ne0\)。对每个素数 \(p\)，把正指数按非降次序排列；令 \(s\) 为各素数指数列表长度的最大值，把较短列表的左端补零至长度 \(s\)。逐位置把相应的 \(p\) 幂相乘，便得到
\[
d_1\mid d_2\mid\cdots\mid d_s.
\]
反过来，任何满足整除链且 \(d_1>1\) 的不变因子分解，其每个素数的指数必非降，去掉左端的零以后恰为该素数的初等因子指数。至少有一个素数整除 \(d_1\)，所以某个列表的长度恰为 \(s\)，不会额外留下全为零的位置。由此不变因子也唯一，\cref{b1:thm:fga-structure}得证。

\ProseParagraphBreak
例如若一个二主群的各层商群阶依次为 \(2^3,2^2,2^2\)，再下一层为一，则指数至少为 \(1,2,3\) 的因子分别有 \(3,2,2\) 个。因此指数一出现一次，指数二不出现，指数三出现两次，该群为 \(C_2\oplus C_8\oplus C_8\)。只知道群的阶为 \(2^7\) 无法作出这一判断，逐层取整数倍才恢复了各个循环因子的长度。
